\medskip

On considère la fonction $f$ définie sur l'intervalle $]2~;~ +\infty[$ par 

\[f(x) = x \ln (x - 2).\]

Une partie de la courbe représentative $\mathcal{C}_f$ de la fonction $f$ est donnée ci-dessous.

\medskip

\begin{center}
\psset{xunit=1cm,yunit=0.2cm,arrowsize=2pt 3}
\begin{pspicture*}(-0.75,-17)(9,17)
\psaxes[linewidth=1.25pt,Dy=5,labelFontSize=\scriptstyle]{->}(0,0)(0,-17)(9,17)
\psplot[plotpoints=2000,linewidth=1.25pt,linecolor=blue]{2.0001}{9}{x 2 sub ln x mul}
\uput[ul](8,13){\blue $\mathcal{C}_f$}
\end{pspicture*}
\end{center}

\medskip

\begin{enumerate}
\item Conjecturer, à l'aide du graphique, le sens de variation de $f$ ses limites aux bornes
de son ensemble de définition ainsi que les éventuelles asymptotes.
\item Résoudre l'équation $f(x) = 0$ sur $]2~;~+\infty[$.
\item Calculer $\begin{gathered}\displaystyle\lim_{\substack{x \to 2\\x > 2}} f(x)\end{gathered}$.

Ce résultat confirme-t-il l'une des conjectures faites à la question 1. ?
\item Démontrer que pour tout $x$ appartenant à $]2~;~+\infty[$ :

\[f'(x) = \ln (x - 2) + \dfrac{x}{x - 2}.\]
\item On considère la fonction $g$ définie sur l'intervalle $]2~;~+\infty[$ par $g(x) = f'(x)$.
	\begin{enumerate}
		\item Démontrer que pour tout $x$ appartenant à $]2~;~+\infty[$, on a:

\[g'(x) = \dfrac{x - 4}{(x - 2)^2}.\]

		\item On admet que $\begin{gathered}\displaystyle\lim_{\substack{x \to 2\\x>2}} g(x) = + \infty\end{gathered}$ et que $\displaystyle\lim_{x \to + \infty} g(x) = +\infty$.

En déduire le tableau des variations de la fonction $g$ sur $]2~;~+\infty[$. On fera apparaître la valeur exacte de l'extremum de la fonction $g$.
		\item En déduire que, pour tout $x$ appartenant à $]2~;~+\infty[$,\: $g(x) > 0$.
		\item En déduire le sens de variation de la fonction $f$ sur $]2~;~+\infty[$.
	\end{enumerate}
\item Étudier la convexité de la fonction $f$ sur $]2~;~+\infty[$ et préciser les coordonnées d'un éventuel point d'inflexion de la courbe représentative de la fonction $f$.
\item Combien de valeurs de $x$ existe-t-il pour lesquelles la courbe représentative de $f$
admet une tangente de coefficient directeur égal à $3$ ?
\end{enumerate}

\bigskip
